\documentclass[12pt,a4paper]{article}

% --- Pakiety podstawowe ---
\usepackage[polish]{babel}
\usepackage[T1]{fontenc}
\usepackage[utf8]{inputenc}
\usepackage{amsmath}
\usepackage{amsfonts}
\usepackage{amssymb}
\usepackage{graphicx}
\usepackage{geometry}
\usepackage{xcolor}
\usepackage{titlesec}
\usepackage{fancyhdr}
\usepackage{longtable}
\usepackage{booktabs}
\usepackage{array}
\usepackage{hyperref}
\usepackage{enumitem}
\usepackage{calc}
% Pandoc-generated macro
\providecommand{\tightlist}{\setlength{\itemsep}{0pt}\setlength{\parskip}{0pt}}

% --- Konfiguracja strony ---
\geometry{margin=2.5cm}
\setlength{\parindent}{0pt}
\setlength{\parskip}{1em}

% --- Kolorystyka ---
\definecolor{NavyBlue}{RGB}{31, 56, 100}
\definecolor{MidBlue}{RGB}{46, 117, 182}
\definecolor{LightBlue}{RGB}{214, 228, 240}
\definecolor{Gray}{RGB}{242, 242, 242}

% --- Stylizacja nagłówków ---
\titleformat{\section}
  {\color{NavyBlue}\normalfont\Large\bfseries}{\thesection}{1em}{}[{\titlerule[1.5pt]}]

\titleformat{\subsection}
  {\color{MidBlue}\normalfont\large\bfseries}{\thesubsection}{1em}{}

% --- Nagłówek i stopka ---
\pagestyle{fancy}
\fancyhf{}
\fancyhead[L]{\small \color{gray} Pakiet Dokumentacji Cyberbezpieczeństwa v1.0}
\fancyhead[R]{\small \color{gray} [NAZWA FIRMY]}
\fancyfoot[C]{\thepage}
\fancyfoot[R]{\small \color{gray} POUFNE}

% --- Ustawienia tabel ---
\setlength{\LTpre}{4pt}
\setlength{\LTpost}{4pt}
\newcolumntype{L}[1]{>{\raggedright\arraybackslash}p{#1}}
\renewcommand{\arraystretch}{1.5}

\begin{document}
\phantomsection\label{doc8}
PROC-008

\section{PROCEDURA SZKOLEŃ I WERYFIKACJI WIEDZY Z
CYBERBEZPIECZEŃSTWA}\label{procedura-szkoleux144-i-weryfikacji-wiedzy-z-cyberbezpieczeux144stwa}

\begin{longtable}[]{@{}>{\raggedright\arraybackslash}p{\dimexpr\linewidth/1000*300-2\tabcolsep\relax}>{\raggedright\arraybackslash}p{\dimexpr\linewidth/1000*700-2\tabcolsep\relax}@{}}
\toprule\noalign{}
\endhead
\bottomrule\noalign{}
\endlastfoot
Wersja: & 1.0 \\
Dotyczy: & Wszyscy pracownicy i współpracownicy \\
Odpowiedzialny: & Administrator IT + Właściciel \\
Częstotliwość: & Szkolenie bazowe: onboarding \textbar{} Odświeżające:
co 12 miesięcy \\
\end{longtable}

\subsection{1. Program szkoleń}\label{program-szkoleux144}

\subsubsection{Szkolenie bazowe (nowi
pracownicy)}\label{szkolenie-bazowe-nowi-pracownicy}

\begin{itemize}
\tightlist
\item
  Czas trwania: min. 2 godziny
\item
  Format: prezentacja + omówienie + quiz
\item
  Zawartość: polityki IT, hasła i MFA, phishing i inżynieria społeczna,
  obsługa incydentów, RODO, bezpieczna praca zdalna
  Dodatkowy moduł dla pracowników księgowości zawierający:
  obsługę faktur, przelewów, deklaracji podatkowych oraz komunikacji z klientami.
\item
  Wymagany wynik quizu: min. \textbf{80\%}
\item
  Powtarzanie przy wyniku poniżej 80\% (max 3 podejścia, potem szkolenie
  indywidualne)
\end{itemize}

\subsubsection{Szkolenie odświeżające (co 12
miesięcy)}\label{szkolenie-odux15bwieux17cajux105ce-co-12-miesiux119cy}

\begin{itemize}
\tightlist
\item
  Czas trwania: min. 1 godzina
\item
  Format: e-learning lub warsztat stacjonarny
\item
  Zawartość: nowe zagrożenia roku, przegląd incydentów z minionego roku,
  nowe procedury.
\item
  Wymagany wynik testu: min. \textbf{80\%}
\end{itemize}

\subsection{2. Testy phishingowe}\label{testy-phishingowe}

Firma przeprowadza symulowane ataki phishingowe w celu weryfikacji
świadomości pracowników:

\begin{longtable}[]{@{}>{\raggedright\arraybackslash}p{\dimexpr\linewidth/1000*270-2\tabcolsep\relax}>{\raggedright\arraybackslash}p{\dimexpr\linewidth/1000*430-2\tabcolsep\relax}>{\raggedright\arraybackslash}p{\dimexpr\linewidth/1000*300-2\tabcolsep\relax}@{}}
\toprule\noalign{}
Parametr & Opis & Wartość docelowa \\
\midrule\noalign{}
\endhead
\bottomrule\noalign{}
\endlastfoot
Częstotliwość & Niezapowiedziane testy & Co kwartał (minimum) \\
Narzędzie & Platforma do symulacji phishingu & GoPhish \\
Próg alertu & Odsetek kliknięć w link & Cel: \textless10\%
pracowników \\
Reakcja na kliknięcie & Natychmiastowe szkolenie remediacyjne &
Obowiązkowe w ciągu 7 dni \\
Raportowanie & Anonimowe wyniki zbiorcze & Co kwartał -- Administrator IT \\
\end{longtable}

\subsection{3. Przegląd Playbooków i
ćwiczenia}\label{przeglux105d-playbookuxf3w-i-ux107wiczenia}

\begin{itemize}
\tightlist
\item
  Pracownicy zapoznają się z Playbookiem Incident Response (PLAY-010) co
  12 miesięcy
\item
  Coroczne ćwiczenie tablicowe (tabletop exercise) -- symulacja
  scenariusza incydentu
\item
  Udokumentowane wyniki ćwiczenia i wnioski (lessons learned)
\end{itemize}

\subsection{4. Rejestr szkoleń (szablon) }\label{rejestr-szkoleux144}

\begin{longtable}[]{@{}>{\raggedright\arraybackslash}p{\dimexpr\linewidth/1000*220-2\tabcolsep\relax}>{\raggedright\arraybackslash}p{\dimexpr\linewidth/1000*200-2\tabcolsep\relax}>{\raggedright\arraybackslash}p{\dimexpr\linewidth/1000*200-2\tabcolsep\relax}>{\raggedright\arraybackslash}p{\dimexpr\linewidth/1000*200-2\tabcolsep\relax}>{\raggedright\arraybackslash}p{\dimexpr\linewidth/1000*180-2\tabcolsep\relax}@{}}
\toprule\noalign{}
Pracownik & Szkolenie & Data & Wynik testu & Certyfikat \\
\midrule\noalign{}
\endhead
\bottomrule\noalign{}
\endlastfoot
\underline{\hspace{3cm}} & Bazowe & \underline{\hspace{2cm}} &
\underline{\hspace{2cm}}\% & {[} {]} TAK \\
\underline{\hspace{3cm}} & Odświeżające & \underline{\hspace{2cm}} &
\underline{\hspace{2cm}}\% & {[} {]} TAK \\
\underline{\hspace{3cm}} & Phishing test & \underline{\hspace{2cm}}
& Kliknął: {[} {]} TAK / NIE & N/A \\
\underline{\hspace{3cm}} & Bazowe & \underline{\hspace{2cm}} &
\underline{\hspace{2cm}}\% & {[} {]} TAK \\
\end{longtable}

\begin{longtable}[]{@{}
  >{\raggedright\arraybackslash}p{\dimexpr\linewidth/1000*500-1\tabcolsep\relax}
  >{\raggedright\arraybackslash}p{\dimexpr\linewidth/1000*500-1\tabcolsep\relax}@{}}
\toprule\noalign{}
\endhead
\bottomrule\noalign{}
\endlastfoot
\begin{minipage}[t]{\linewidth}\raggedright
Zatwierdził:

Podpis: \underline{\hspace{3cm}}

Data: \underline{\hspace{4cm}}
\end{minipage} & \begin{minipage}[t]{\linewidth}\raggedright
Stanowisko / Tytuł:

Imię i nazwisko: \underline{\hspace{3cm}}

Funkcja: \underline{\hspace{3cm}}
\end{minipage} \\
\end{longtable}

\phantomsection\label{doc10}
PLAY-010

\section{PLAYBOOK SYTUACJI NIETYPOWYCH -- INCIDENT
RESPONSE}\label{playbook-sytuacji-nietypowych-incident-response}

\begin{longtable}[]{@{}>{\raggedright\arraybackslash}p{\dimexpr\linewidth/1000*300-2\tabcolsep\relax}>{\raggedright\arraybackslash}p{\dimexpr\linewidth/1000*700-2\tabcolsep\relax}@{}}
\toprule\noalign{}
\endhead
\bottomrule\noalign{}
\endlastfoot
Wersja: & 1.0 \\
Dotyczy: & Wszyscy pracownicy \\
Odpowiedzialny: & Właściciel / Administrator IT \\
Aktualizacja: & Po każdym incydencie i co 12 miesięcy \\
\end{longtable}

\textbf{ALARM:} Ten dokument musi być dostępny OFFLINE (wydrukowany) w widocznym
miejscu. W przypadku incydentu sieć może być niedostępna.

\subsection{1. Kontakty alarmowe -- WYPEŁNIĆ I
WYDRUKOWAĆ}\label{kontakty-alarmowe-wypeux142niux107-i-wydrukowaux107}

\begin{longtable}[]{@{}>{\raggedright\arraybackslash}p{\dimexpr\linewidth/1000*270-2\tabcolsep\relax}>{\raggedright\arraybackslash}p{\dimexpr\linewidth/1000*250-2\tabcolsep\relax}>{\raggedright\arraybackslash}p{\dimexpr\linewidth/1000*250-2\tabcolsep\relax}>{\raggedright\arraybackslash}p{\dimexpr\linewidth/1000*230-2\tabcolsep\relax}@{}}
\toprule\noalign{}
Rola & Imię i Nazwisko & Telefon / E-mail & Zakres odpowiedzialności \\
\midrule\noalign{}
\endhead
\bottomrule\noalign{}
\endlastfoot
Właściciel firmy (CISO) & \underline{\hspace{3cm}} &
\underline{\hspace{3cm}} & Główna decyzja, autoryzacja działań \\
Administrator IT & \underline{\hspace{3cm}} &
\underline{\hspace{3cm}} & Techniczne działanie i izolacja \\
Osoba do kontaktu z Policją & \underline{\hspace{3cm}} &
\underline{\hspace{3cm}} & Zgłoszenia prawne, przestępstwa \\
Osoba do kontaktu z mediami & \underline{\hspace{3cm}} &
\underline{\hspace{3cm}} & JEDEN głos firmy, komunikacja zewn. \\
Osoba do kontaktu z UODO & \underline{\hspace{3cm}} &
\underline{\hspace{3cm}} & Naruszenia RODO (72h) \\
Prawnik firmy & \underline{\hspace{3cm}} &
\underline{\hspace{3cm}} & Porady prawne, kontakty z organami \\
Ubezpieczyciel (cyber) & \underline{\hspace{3cm}} &
\underline{\hspace{3cm}} & Zgłoszenie szkody \\
Zewnętrzna firma IT (backup) & \underline{\hspace{3cm}} &
\underline{\hspace{3cm}} & Wsparcie techniczne awaryjne \\
\end{longtable}

Ważne numery: Policja: 997 \textbar{} CERT Polska: cert.pl \textbar{}
UODO: uodo.gov.pl \textbar{} tel. UODO: 22 531 03 00

\subsection{2. Fazy reagowania na
incydent}\label{fazy-reagowania-na-incydent}

\subsubsection{FAZA 1: Identyfikacja (0--30
min)}\label{faza-1-identyfikacja-030-min}

\begin{enumerate}
\tightlist
\item
  Pracownik wykrywa anomalię → \textbf{NATYCHMIASTOWO} informuje
  administratora IT i właściciela
\item
  Administrator IT ocenia powagę incydentu wg klasyfikacji poniżej
\item
  Rejestracja incydentu w Rejestrze Incydentów (godzina, opis, reporter)
\end{enumerate}

\subsubsection{FAZA 2: Izolacja (30--60
min)}\label{faza-2-izolacja-3060-min}

\begin{enumerate}
\tightlist
\item
  Odizolowanie dotkniętych systemów od sieci (jeśli dotyczy) -- kabel
  LAN + WiFi
\item
  Zmiana haseł do zagrożonych kont
\item
  Zabezpieczenie dowodów (logi, kopie plików) -- \textbf{NIE USUWAĆ
  niczego}
\end{enumerate}

\subsubsection{FAZA 3: Ograniczenie i
likwidacja}\label{faza-3-ograniczenie-i-likwidacja}

\begin{enumerate}
\tightlist
\item
  Blokada ruchu sieciowego do/ze źródeł ataku na firewall
\item
  Przywracanie z kopii zapasowej (wg DRP-011)
\item
  Identyfikacja wektora ataku i usunięcie złośliwego kodu
\end{enumerate}

\subsection{3. Matryca powiadamiania}\label{matryca-powiadamiania}

\begin{longtable}[]{@{}>{\raggedright\arraybackslash}p{\dimexpr\linewidth/1000*270-2\tabcolsep\relax}>{\raggedright\arraybackslash}p{\dimexpr\linewidth/1000*250-2\tabcolsep\relax}>{\raggedright\arraybackslash}p{\dimexpr\linewidth/1000*250-2\tabcolsep\relax}>{\raggedright\arraybackslash}p{\dimexpr\linewidth/1000*230-2\tabcolsep\relax}@{}}
\toprule\noalign{}
Kogo powiadomić & Kiedy & Kto powiadamia & Termin \\
\midrule\noalign{}
\endhead
\bottomrule\noalign{}
\endlastfoot
UODO (Urząd Ochrony Danych) & Naruszenie danych osobowych z ryzykiem &
Właściciel / IOD & \textbf{72 godziny} \\
Klienci poszkodowani & Wysokie ryzyko dla ich danych & Właściciel & Bez
zwłoki \\
Policja / CERT Polska & Przestępstwo, atak hakerski, ransomware &
Wyznaczona osoba & Niezwłocznie \\
Media / PR & Jeśli incydent jest/może być publiczny & Właściciel (JEDEN
głos) & Po konsultacji z prawnikiem \\
Ubezpieczyciel & Szkoda objęta polisą cyber & Właściciel & 24--48
godzin \\
\end{longtable}

\subsection{4. Scenariusze -- reakcje}\label{scenariusze-reakcje}

\subsubsection{Scenariusz A: Phishing / kompromitacja
konta}\label{scenariusz-a-phishing-kompromitacja-konta}

\begin{itemize}
\tightlist
\item
  Natychmiast zmień hasło i zresetuj MFA dla konta
\item
  Przegląd aktywności konta (logi) za ostatnie 30 dni
\item
  Sprawdź czy dane były eksportowane lub przesyłane
\item
  Poinformuj pozostałych pracowników (alert phishingowy)
\end{itemize}

\subsubsection{Scenariusz B:
Ransomware}\label{scenariusz-b-ransomware}

\begin{itemize}
\tightlist
\item
  \textbf{NATYCHMIASTOWO} odłącz zainfekowane urządzenie od sieci (kabel
  + WiFi)
\item
  \textbf{NIE PŁACIĆ OKUPU} bez konsultacji z prawnikiem i
  ubezpieczycielem
\item
  Uruchom procedurę DRP-011
\item
  Zgłoś do Policji i CERT Polska (cert.pl)
\item
  Zidentyfikuj wektor ataku i załataj podatność
\end{itemize}

\subsubsection{Scenariusz C: Wyciek danych
klientów}\label{scenariusz-c-wyciek-danych-klientuxf3w}

\begin{itemize}
\tightlist
\item
  Zidentyfikuj jakie dane wyciekły i ile osób dotyczy
\item
  Ocena ryzyka wg RODO -- czy zgłosić do UODO (72h)?
\item
  Powiadom poszkodowanych klientów
\item
  Zabezpiecz dowody (nie usuwaj logów)
\end{itemize}

\subsubsection{Scenariusz D: Fraud finansowy / zmiana rachunku}\label{scenariusz-D-kompromitacja-systemu}
\begin{itemize}
\tightlist
\item
\textbf{NATYCHMIASTOWO} wstrzymaj przelew
\item
Skontaktować się z bankiem
\item
Skontaktowanie się z klientem
\item
Zabezpiecz korespondencję e-mail
\item
Zgłoś incydent do Policji
\item 
Zarejestruj incydent w Rejestrze Incydentów
\end{itemize}

\begin{longtable}[]{@{}>{\raggedright\arraybackslash}p{\dimexpr\linewidth/1000*220-2\tabcolsep\relax}>{\raggedright\arraybackslash}p{\dimexpr\linewidth/1000*200-2\tabcolsep\relax}>{\raggedright\arraybackslash}p{\dimexpr\linewidth/1000*200-2\tabcolsep\relax}>{\raggedright\arraybackslash}p{\dimexpr\linewidth/1000*200-2\tabcolsep\relax}>{\raggedright\arraybackslash}p{\dimexpr\linewidth/1000*180-2\tabcolsep\relax}@{}}
\toprule\noalign{}
Data/Czas & Opis incydentu & Działania podjęte & Status & Zamknięty \\
\midrule\noalign{}
\endhead
\bottomrule\noalign{}
\endlastfoot
\underline{\hspace{2cm}} &
\underline{\hspace{4cm}} &
\underline{\hspace{4cm}} & {[} {]} Otwarty &
\underline{\hspace{2cm}} \\
\underline{\hspace{2cm}} &
\underline{\hspace{4cm}} &
\underline{\hspace{4cm}} & {[} {]} Otwarty &
\underline{\hspace{2cm}} \\
\end{longtable}

\begin{longtable}[]{@{}
  >{\raggedright\arraybackslash}p{\dimexpr\linewidth/1000*500-1\tabcolsep\relax}
  >{\raggedright\arraybackslash}p{\dimexpr\linewidth/1000*500-1\tabcolsep\relax}@{}}
\toprule\noalign{}
\endhead
\bottomrule\noalign{}
\endlastfoot
\begin{minipage}[t]{\linewidth}\raggedright
Zatwierdził:

Podpis: \underline{\hspace{3cm}}

Data: \underline{\hspace{4cm}}
\end{minipage} & \begin{minipage}[t]{\linewidth}\raggedright
Stanowisko / Tytuł:

Imię i nazwisko: \underline{\hspace{3cm}}

Funkcja: \underline{\hspace{3cm}}
\end{minipage} \\
\end{longtable}
\end{document}